\exercisegroup

\begin{enumerate}
\def\labelenumi{\arabic{enumi})}
\tightlist
\item
  设 \(A\) 是一个无穷集。
\end{enumerate}

\begin{enumerate}
\def\labelenumi{\alph{enumi})}
\tightlist
\item
  设 E 是 \(\overline{\mathcal{K}(\mathbf{A})}\) 在 \(\mathbf{R}\) 上的 Banach 空间，由满足下列条件的实数族 \(x = (x_{\alpha})_{\alpha \in \mathbf{A}}\) 组成：\(\alpha \mapsto x_{\alpha}\) 对于 A 的有限子集的余集滤子趋于 0，赋以范数 \(\| x\| = \sup_{\alpha \in \mathbf{A}}|x_{\alpha}|\)（当 \(\mathbf{A} = \mathbf{N}\) 时，这个空间记作 \(c_0\) 或 \(c_0(\mathbf{N})\)。证明 E 上的每个连续线性型都可以唯一地写成 \(x \mapsto \sum_{\alpha \in \mathbf{A}} u_{\alpha} x_{\alpha}\)，其中 \((u_{\alpha})_{\alpha \in \mathbf{A}}\) 是一个满足 \(\sum_{\alpha \in \mathbf{A}} |u_{\alpha}| < +\infty\) 的族；于是 E 的对偶 \(\mathrm{E}'\)
\end{enumerate}

可以（作为非拓扑向量空间）与空间 \(\ell^1 (\mathbf{A})\) 视为同一 (I, 第 4 页, 例)。\(b)\) 设 \(\mathrm{F}\) 是 Banach 空间 \(\ell^1 (\mathbf{A})\) (I, 第 4 页, 例)（当 \(\mathbf{A} = \mathbf{N}\) 时，也记作 \(\ell^1\)）。证明 \(\mathrm{F}\) 上的每个连续线性型都可以唯一地写成 \(x\mapsto \sum_{\alpha \in \mathbf{A}}u_{\alpha}x_{\alpha}\)，其中 \((u_{\alpha})_{\alpha \in \mathbf{A}}\) 是一个有界的实数族；于是 \(\mathrm{F}'\)（\(\mathrm{F}\) 的对偶）可以（作为非拓扑向量空间）与空间 \(\mathcal{B}(\mathbf{A}) = \ell^{\infty}(\mathbf{A})\) 视为同一。

\begin{enumerate}
\def\labelenumi{\alph{enumi})}
\setcounter{enumi}{2}
\tightlist
\item
  设 B 是任意一个集，\((c_{\alpha \beta})_{(\alpha ,\beta)\in \mathrm{A}\times \mathrm{B}}\) 是任意一个 \(\geqslant 0\) 的数族。设 G 是由满足下列条件的实数族 \(x = (x_{\alpha})_{\alpha \in \mathbf{A}}\) 所成的向量空间：对每个 \(\beta \in \mathbf{B}\)，我们有 \(p_{\beta}(x) = \sum_{\alpha \in \mathbf{A}}c_{\alpha \beta}|x_{\alpha}| < + \infty\)；诸 \(p_{\beta}\) 是 G 上的半范数。为使 G（赋以由该半范数族定义的拓扑）是 Hausdorff 的，必须且只需：对每个 \(\alpha \in \mathbf{A}\)，至少存在一个 \(\beta \in \mathbf{B}\) 使得 \(c_{\alpha \beta} > 0\)。证明这时 \(\mathbf{G}\) 完备，且 \(\mathbf{G}\) 上的每个连续线性型都可以唯一地写成 \(x \mapsto \sum_{\alpha \in \mathbf{A}} u_{\alpha} x_{\alpha}\)，其中
\end{enumerate}

\((u_{\alpha})_{\alpha \in \mathbf{A}}\) 是满足下列条件的实数族：存在有限多个指标 \(\beta_{i}\in \mathrm{B}(1\leqslant i\leqslant n)\) 与一个数 \(a > 0\)，使得 \(|u_{\alpha}|\leqslant a.c_{\alpha \beta_i}\) 对所有 \(\alpha \in \mathbf{A}\) 与 \(1\leqslant i\leqslant n\) 成立。证明其逆，并把这些结果推广到标量域为 \(\mathbf{C}\) 的情形。

\begin{enumerate}
\def\labelenumi{\arabic{enumi})}
\setcounter{enumi}{1}
\tightlist
\item
  \begin{enumerate}
  \def\labelenumii{\alph{enumii})}
  \tightlist
  \item
    设 F 与 G 是处于分离对偶中的两个向量空间。证明：若 F 对于 \(\sigma(\mathrm{F}, \mathrm{G})\) 相对有界 (III, 第 43 页, 习题 6)，则 G 对于 \(\sigma(\mathrm{G}, \mathrm{F})\) 相对有界。
  \end{enumerate}
\end{enumerate}

\begin{enumerate}
\def\labelenumi{\alph{enumi})}
\setcounter{enumi}{1}
\item
  设 \(\mathbf{F}\) 是一个向量空间，\(\mathbf{G}_1, \mathbf{G}_2\) 是 \(\mathbf{F}^*\) 的两个向量子空间，使得 \(\mathbf{F}\) 与 \(\mathbf{G}_1\) 以及与 \(\mathbf{G}_2\) 都处于分离对偶中。证明：若 \(\mathbf{F}\) 对于 \(\sigma(\mathbf{F}, \mathbf{G}_1)\) 与 \(\sigma(\mathbf{F}, \mathbf{G}_2)\) 都相对有界，则它对于 \(\sigma(\mathbf{F}, \mathbf{G}_1 + \mathbf{G}_2)\) 也相对有界。
\item
  假设 F 具有可数基。证明：对于 \(\mathbf{F}^*\) 的每个与 F 处于分离对偶且具有可数基的向量子空间 G，F 对于 \(\sigma(\mathrm{F},\mathrm{G})\) 相对有界（用归纳法定义 F 与 G 的两个基 \((a_n)\)、\((b_n)\)，使得 \(\langle a_m, b_n \rangle = \delta_{mn}\)）。
\end{enumerate}

\begin{enumerate}
\def\labelenumi{\arabic{enumi})}
\setcounter{enumi}{2}
\tightlist
\item
  设 F 是一个向量空间。证明：对于拓扑 \(\sigma(F, F^{*})\)，F 的每个有界子集都是有限维的。由此推出，若 F 是无穷维的，则在 F 的完备化 \(\hat{F}\)（对于 \(\sigma(F, F^{*})\)）中，存在不含于 F 的任何有界子集的闭包中的紧子集（参看 II, 第 52 页, 命题 10）。
\end{enumerate}

¶ 4) 设 F 与 G 是处于分离对偶中的两个向量空间，G（相应地，F）当后者赋以拓扑 σ(F, G)（相应地，σ(G, F)）时与 F（相应地，G）的对偶视为同一。设 S（相应地，T）是 F（相应地，G）的一个覆盖，由对于 σ(F, G)（相应地，σ(G, F)）凸、均衡且有界的子集组成。证明下列命题等价：

\(\alpha)\) 每个集 \(\mathbf{M} \in \mathfrak{S}\) 对于 \(\mathfrak{T}\)-拓扑预紧。

β) 每个集 \(\mathbf{N} \in \mathfrak{T}\) 对于 \(\mathfrak{S}\)-拓扑预紧。

\(\gamma)\) 在每个集 \(\mathbf{M} \in \mathfrak{S}\) 上，由 \(\mathfrak{T}\)-拓扑诱导的拓扑与由 \(\sigma(\mathrm{F}, \mathrm{G})\) 诱导的拓扑相同。

\(\delta)\) 在每个集 \(\mathbf{N} \in \mathfrak{T}\) 上，由 \(\mathfrak{S}\)-拓扑诱导的拓扑与由 \(\sigma(\mathbf{G}, \mathbf{F})\) 诱导的拓扑相同。

（利用 III, 第 17 页 的命题 5 证明 \(\alpha\) ) 蕴含 \(\delta\) )，并利用 II, 第 74 页 的习题 1 证明 \(\delta\) ) 蕴含 \(\beta\)。）

\begin{enumerate}
\def\labelenumi{\arabic{enumi})}
\setcounter{enumi}{4}
\item
  设 E 与 F 是两个局部凸 Hausdorff 空间，S 是 E 的子集的一个族。为使空间 \(\mathcal{L}(\mathrm{E};\mathrm{F})\) 上的 S-拓扑与向量空间结构相容，必须（且只需，参看 III, 第 13 页, 推论）S 的每个集都在 E 中有界。
\item
  \begin{enumerate}
  \def\labelenumii{\alph{enumii})}
  \tightlist
  \item
    设 E 是一个局部凸 Hausdorff 空间。对 E 上的每个超滤子 U，设 \(U'\) 是以 U 的集的凸包为基本邻域系的滤子。证明：为使 E 的一点对于弱化拓扑是 U 的极限，必须且只需它对于初始拓扑是 \(U'\) 的接触点（利用 II, 第 84 页 的习题 11）。
  \end{enumerate}
\end{enumerate}

\begin{enumerate}
\def\labelenumi{\alph{enumi})}
\setcounter{enumi}{1}
\tightlist
\item
  设 A 是向量空间 E 的一个凸子集，\(T_{1}, T_{2}\) 是 E 上两个局部凸 Hausdorff 拓扑，\(T_{1}^{\prime}, T_{2}^{\prime}\) 是相应的弱化拓扑。证明：若 \(T_{1}\) 在 A 上诱导的拓扑比 \(T_{2}\) 诱导的拓扑细，则 \(T_{1}^{\prime}\) 在 A 上诱导的拓扑比 \(T_{2}^{\prime}\) 在 A 上诱导的拓扑细。
\end{enumerate}

¶ 7) 设 \(\mathbf{F}\) 是直和空间 \(\mathbf{R}^{(\mathbf{N})}\)，\(\mathbf{G}\) 是空间 \(\ell^1 (\mathbf{N})\) (I, 第 4 页, 例)；\(\mathbf{F}\) 与 \(\mathbf{G}\) 通过双线性型

\[
(x, y) \mapsto \sum_ {n} \xi_ {n} \eta_ {n}
\]

置于对偶中，其中 \(x = (\xi_n)\in \mathbf{F}\)，\(y = (\eta_{n})\in \mathbf{G}\)。

\begin{enumerate}
\def\labelenumi{\alph{enumi})}
\tightlist
\item
  证明：在 F 中，每个凸且对于 \(\sigma(F, G)\) 紧的集 K 都是有限维的。（假设相反；设 \((n_k)\) 是一个严格递增的整数序列 \(>0\)，\((a_k)\) 是 K 的点的序列，使得 \(a_k\) 的指标 \(>n_k\) 的分量均为零，但指标 \(n_k\) 的分量 \(\neq 0\)。证明存在一个序列 \((t_k)\)（其中的数 \(>0\)），使得 \(\sum_{k} t_k < +\infty\)，并且在由实数的所有有界序列组成的 Banach 空间 \(\mathcal{B}(N)\) 中，
\end{enumerate}

点 \(\sum_{k} t_k a_k\) 的指标 \(n_i\) 的分量对所有的 \(i\) 都非零；由此推出部分和序列

\(s_{p} = \sum_{k=1}^{p} t_k a_k\) 在 F 中对于 \(\sigma(F, G)\) 不可能有极限点。

\begin{enumerate}
\def\labelenumi{\alph{enumi})}
\setcounter{enumi}{1}
\item
  证明：在 F 中存在对于 \(\sigma(\mathrm{F},\mathrm{G})\) 的无穷维紧集（注意 G 是 F 对于由赋范空间 \(\mathcal{B}(\mathbf{N})\) 的拓扑所诱导的拓扑的对偶）。证明在 F 中存在对于 \(\sigma(\mathrm{F},\mathrm{G})\) 的预紧集，它们不是相对紧的。
\item
  由 a) 与 b) 推出 \(\tau(\mathbf{G},\mathbf{F})=\sigma(\mathbf{G},\mathbf{F})\)，但 \(\tau(\mathbf{G},\mathbf{F})\) 不同于在 F 的对于 \(\sigma(\mathbf{F},\mathbf{G})\) 紧的子集上的一致收敛拓扑。
\end{enumerate}

\begin{enumerate}
\def\labelenumi{\arabic{enumi})}
\setcounter{enumi}{7}
\tightlist
\item
  设 E 是一个无穷维的局部凸可度量化空间，\(E'\) 是它的对偶。为使拓扑 \(\tau(\mathrm{E}, \mathrm{E}')\) 与弱化拓扑 \(\sigma(\mathrm{E}, \mathrm{E}')\) 相同，必须且只需 E 同构于乘积空间 \(R^{N}\)（相应地，\(C^{N}\)）的一个处处稠密的子空间。（注意在 \(E'\) 中，对于由等度连续集组成的有界结构，存在一个由有限维集组成的可数基（III, 第 1 页）；由此推出向量空间 \(E'\) 具有可数基。）
\end{enumerate}

给出一个局部凸 Hausdorff 空间 \(\mathbf{E}\) 的例子，使得初始拓扑与拓扑 \(\sigma(\mathrm{E},\mathrm{E}')\) 和 \(\tau(\mathrm{E},\mathrm{E}')\) 互不相同。

\begin{enumerate}
\def\labelenumi{\arabic{enumi})}
\setcounter{enumi}{8}
\tightlist
\item
  \begin{enumerate}
  \def\labelenumii{\alph{enumii})}
  \tightlist
  \item
    设 E 是一个局部凸、Hausdorff、有界型且拟完备的空间。为使拓扑 \(\tau(\mathrm{E}^{\prime},\mathrm{E})\) 与 \(\sigma(\mathrm{E}^{\prime},\mathrm{E})\) 在 E 的对偶 \(E^{\prime}\) 上相同，必须且只需 E 的拓扑是最细的局部凸拓扑（证明 E 的每个有界子集都是有限维的）。
  \end{enumerate}
\end{enumerate}

\begin{enumerate}
\def\labelenumi{\alph{enumi})}
\setcounter{enumi}{1}
\tightlist
\item
  设 E 是一个局部凸 Hausdorff 空间，\(E'\) 是它的对偶。为使强拓扑 \(\beta(E', E)\)（在 \(E'\) 上）与弱拓扑 \(\sigma(E', E)\) 相同，必须且只需与 E 相伴的有界型空间（III, 第 40 页, 习题 1）的拓扑是 E 上最细的局部凸拓扑。
\end{enumerate}

\begin{enumerate}
\def\labelenumi{\arabic{enumi})}
\setcounter{enumi}{9}
\tightlist
\item
  设 \(\mathbf{E}\) 是一个局部凸 Hausdorff 空间，\(\mathbf{E}'\) 是它的对偶。证明下列命题等价：
\end{enumerate}

α) E 是桶型空间；

β) E' 的每个弱有界子集都等度连续；

\(\gamma)\) \(\mathbf{E}'\) 的每个弱有界子集都是相对弱紧的，且 \(\mathbf{E}\) 的拓扑是 \(\tau (\mathbf{E},\mathbf{E}')\)。

\begin{enumerate}
\def\labelenumi{\arabic{enumi})}
\setcounter{enumi}{10}
\item
  设 E 是一个局部凸 Hausdorff 空间，\(E'\) 是它的对偶，S 是由 E 的有界子集组成的一个覆盖；我们给 \(E'\) 赋以 S-拓扑。证明：为使双线性型 \((x, x') \mapsto \langle x, x' \rangle\) 在 \(E \times E'\) 上连续，必须且只需 E 的拓扑可由单个范数定义，且 S-拓扑是 \(E'\) 上的强拓扑（参看 III, 第 37 页, 习题 2）。
\item
  设 \(E\) 是一个局部凸 Hausdorff 空间，\(E'\) 是它的对偶。
\end{enumerate}

\begin{enumerate}
\def\labelenumi{\alph{enumi})}
\item
  为使在 \(E'\) 中存在一个弱有界且吸收的集，必须且只需 E 的拓扑比一个赋范空间的拓扑粗（参看 IV, 第 47 页, 习题 2）。
\item
  为使在 \(E'\) 中存在一个等度连续且弱完全的集，必须且只需 E 的拓扑比一个赋范空间的拓扑细。
\item
  为使在 \(E'\) 中存在一个等度连续且吸收的集，必须且只需 E 的拓扑可由单个范数定义。
\end{enumerate}

\begin{enumerate}
\def\labelenumi{\arabic{enumi})}
\setcounter{enumi}{12}
\tightlist
\item
  设 \(\mathbf{F}\) 与 \(\mathbf{G}\) 是 \(\mathbf{R}\) 上处于分离对偶中的两个向量空间。
\end{enumerate}

\begin{enumerate}
\def\labelenumi{\alph{enumi})}
\item
  设 A 是 F 中的一个凸集，不含原点，且对于拓扑 \(\sigma(\mathbf{F}, \mathbf{G})\) 紧；设 C 是以 0 为顶点且由 A 生成的凸锥。证明 G 中的极锥 \(C^{\circ}\) 对于拓扑 \(\tau(\mathbf{G}, \mathbf{F})\) 有一个内点：
\item
  反之，设 C 是一个以 0 为顶点的真凸锥，对于 \(\sigma(F, G)\) 闭，且对于拓扑 \(\tau(F, G)\) 有一个内点。证明在 G 中存在一个超平面 H，它对于 \(\sigma(G, F)\) 闭、不含原点，使得 \(H \cap C^{\circ}\) 对于 \(\sigma(G, F)\) 紧，且 \((H \cap C^{\circ}) \cup \{0\}\) 生成 \(C^{\circ}\)。
\end{enumerate}

\begin{enumerate}
\def\labelenumi{\arabic{enumi})}
\setcounter{enumi}{13}
\tightlist
\item
  设 E 是一个局部凸 Hausdorff 且拟完备的空间，\(E'\) 是它的对偶。证明：在 \(E_{1}'\) 上，紧收敛拓扑是与 E 和 \(E'\) 之间的对偶相容、且与 \(\sigma(E', E)\) 在 \(E'\) 的每个等度连续子集上诱导相同拓扑的诸拓扑中最细的一个（参看 IV, 第 48 页, 习题 4）。
\end{enumerate}

¶ 15) a) 设 E 是一个局部凸 Hausdorff 空间，A 是 E 中的一个凸且均衡的集，u 是 E 上的一个线性型。证明：若 \(\mathrm{A} \cap u^{-1}(0)\) 相对于 A 是闭的，则 u 在 A 上的限制连续。（若不然，证明 0 将属于 A 与一个超平面 \(u^{-1}(\alpha)\)（\(\alpha \neq 0\)）的交的闭包；由此推出，若 \(b \in A\) 满足 \(u(b) = -\alpha\)，则点 \(\frac{1}{2}b\) 将属于 \(\mathrm{A} \cap u^{-1}(0)\) 的闭包。）

\begin{enumerate}
\def\labelenumi{\alph{enumi})}
\setcounter{enumi}{1}
\item
  设 E 是一个实的局部凸 Hausdorff 且完备的空间，\(E'\) 是它的对偶。证明：若超平面 \(H'\)（\(E'\) 的）与每个等度连续且弱闭的子集 \(M' \subset E'\) 的交都是弱闭的，则 \(H'\) 是弱闭的（利用 a) 与 III, 第 21 页, 推论 1）。
\item
  设 E 是一个局部凸 Hausdorff 空间，\(E'\) 是它的对偶，C 是 E 中的一个凸、均衡且闭的集。设 u 是 E 上的一个线性型；证明：若 u 在 C 上的限制对于初始拓扑连续，则它对于 \(\sigma(\mathrm{E}, \mathrm{E}')\) 也连续（利用 a)）。用一个例子证明 u 在由 C 生成的向量子空间 M 上的限制未必连续（取 \(E = R^{(N)}\)，赋以范数 \(\|x\| = \sup_{n} |\xi_{n}|\)，并对 C 取一个生成 E 的适当凸集）。
\end{enumerate}

\begin{enumerate}
\def\labelenumi{\arabic{enumi})}
\setcounter{enumi}{15}
\tightlist
\item
  设 F 与 G 是处于分离对偶中的两个向量空间；把 F 上所有线性型 \(x'\)（在 F 的每个对于 \(\sigma(F, G)\) 有界的子集上都有界者）所成的集，称为 G 在 F 的代数对偶 \(F^{*}\) 中的包封；记这个向量子空间为 \(\tilde{G}\)，它含于 \(F^{*}\)，且当 F 赋以与 F 和 G 之间的对偶相容的诸拓扑之一的相配有界型空间（III, 第 40 页, 习题 1）的拓扑时，它就是 F 的对偶。称 G 在 \(F^{*}\) 中是包封的，如果 \(\tilde{G} = G\)。
\end{enumerate}

\begin{enumerate}
\def\labelenumi{\alph{enumi})}
\tightlist
\item
  设 M 是 F 的一个对于拓扑 \(\sigma(\mathrm{F},\mathrm{G})\) 闭的向量子空间。证明：若 G 在 \(F^{*}\) 中是包封的，且 F/M 赋以拓扑 \(\sigma(\mathrm{F}/\mathrm{M},\mathrm{M}^{\circ})\)，则它的对偶在 \((\mathrm{F}/\mathrm{M})^{*}\) 中是包封的。
\item
  设 E 是一个 Hausdorff 局部凸空间，\(E'\) 是它的对偶。为使 E 是有界型的，必须且只需它的拓扑与 \(\tau(\mathrm{E},\mathrm{E}')\) 相同，且 \(E'\) 在 \(E^{*}\) 中是包封的。
\item
  设 \((\mathrm{E}_{i})_{i\in\mathbb{I}}\) 是一个 Hausdorff 局部凸空间族，F 是 \(E_{i}\) 的直和空间，赋以 I, 第 24 页, 习题 14 中定义的拓扑。证明 F 的对偶典范同构于乘积 \(\prod_{i\in I}E_{i}'\)（诸 \(E_{i}\) 的对偶的乘积）中由
\end{enumerate}

所有族 \((x_{i}^{\prime})\)（除可数多个指标外 \(x_{i}^{\prime}=0\)）组成的子空间。（设 \(V_{i}\) 是 \(E_{i}\) 中 0 的任意邻域。证明：若对一个不可数的指标集有 \(x_{i}^{\prime}\neq0\)，则存在一个数 \(\alpha>0\) 与一个不可数集 H⊂I，使得存在 \(x_{i}\in V_{i}\)，它满足 \(\langle x_{i},x_{i}^{\prime}\rangle\geqslant\alpha\)（对所有 \(i\in H\)）；由此推出 \((x_{i}^{\prime})\) 不可能属于 F'。）

\begin{enumerate}
\def\labelenumi{\alph{enumi})}
\setcounter{enumi}{3}
\tightlist
\item
  证明：若 I 不可数，则 \(\mathbf{F}'\) 在 \(\mathbf{F}^*\) 中不是包封的；若取 \(\mathrm{E}_i = \mathbf{R}\)（对所有 \(i \in \mathrm{I}\)），则赋以强拓扑的 \(\mathbf{F}'\) 不完备，且在 \(\mathbf{F}'\) 中存在不是相对弱紧的强有界子集。
\end{enumerate}

\begin{enumerate}
\def\labelenumi{\arabic{enumi})}
\setcounter{enumi}{16}
\item
  设 E 是一个 Hausdorff 局部凸空间，\(E'\) 是它的对偶；在 \(E'\) 中，设 \(B_{1}\) 是凸等度连续集的族，\(B_{2}\) 是凸的相对弱紧子集的族，\(B_{3}\) 是凸的强有界子集的族，\(B_{4}\) 是凸的弱有界子集的族。于是 \(B_{1} \subset B_{2} \subset B_{3} \subset B_{4}\)。给出一个使这四个集族互不相同的空间 E 的例子（对 E 取三个空间的乘积，对它们分别有 \(B_{1} \neq B_{2}\)（参看 IV, 第 49 页, 习题 8），\(B_{2} \neq B_{3}\)（习题 16, d)），\(B_{3} \neq B_{4}\)（III, 第 23 页, 注 2）。
\item
  设 E、F 是两个向量空间，\(E^{*}\)、\(F^{*}\) 分别是它们的代数对偶。证明：若 u 是一个从 \(F^{*}\) 到 E 中的线性映射，它对于拓扑 \(\sigma(F^{*}, F)\) 与 \(\sigma(E, E^{*})\) 连续，则 \(u(F^{*})\) 是有限维的。（利用 IV, 第 27 页 的命题 2 证明 u 是一个严格态射，并由此推出 \(u(F^{*})\) 是 E 的一个极小型子空间（II, 第 85 页, 习题 13）；考虑这个子空间的有界集以结束证明。）
\item
  设 E 与 F 是两个 Hausdorff 局部凸空间，\(E'\) 与 \(F'\) 是它们的对偶。对从 E 到 F 中的连续线性映射所成的空间 \(\mathcal{L}(\mathrm{E};\mathrm{F})\) 的每个子集 H，用 \({}^{t}H\) 表示映射 \(u\in H\) 的转置所成的集。对 E 的每个子集 M（相应地，子集 \(N'\)，它属于 \(F'\)），用 \(\mathrm{H}(\mathrm{M})\)（相应地，\({}^{t}\mathrm{H}(\mathrm{N}')\)）表示当 u 遍历 H 时诸集 \(u(\mathrm{M})\)（相应地，\({}^{t}u(\mathrm{N}')\)）的并。a) 为使 H 等度连续，必须且只需：对每个等度连续子集 \(N'\)（\(F'\) 的），\({}^{t}\mathrm{H}(\mathrm{N}')\) 是 \(E'\) 的一个等度连续子集。
\end{enumerate}

\begin{enumerate}
\def\labelenumi{\alph{enumi})}
\setcounter{enumi}{1}
\item
  设 \(\mathfrak{S}\) 是 E 的有界子集的一个集。证明：H 在 \(\mathcal{L}(\mathrm{E};\mathrm{F})\) 中对于 \(\mathfrak{S}\) -拓扑有界，当且仅当对每个 \(y' \in \mathrm{F}'\)，\({}^{t}\mathrm{H}(y')\) 在 \(\mathrm{E}'\) 中对于 \(\mathfrak{S}\) -拓扑有界。
\item
  设 S 是 E 的有界子集的一个集，J 是 F 的有界子集的一个集，构成 F 的一个适配有界结构（III, 第 3 页, 定义 4）。假设 E' 赋以 S-拓扑，F' 赋以 J-拓扑。于是 \({}^{t}\mathrm{H}\) 等度连续当且仅当对每个集 B ∈ S，H(B) 属于 J。特别地，\({}^{t}\mathrm{H}\) 对于 F' 与 E' 上的强拓扑等度连续，当且仅当 H 在 L(E; F) 中对于有界收敛拓扑有界。
\end{enumerate}

\(d\) ) 由 \(b\) ) 与 \(c\) ) 推出：若 \({}^t\mathrm{H}\) 在 \(\mathcal{L}(\mathrm{F}';\mathrm{E}')\) 中对于简单收敛拓扑有界（当 \(\mathrm{F}'\) 与 \(\mathrm{E}'\) 赋以强拓扑时），则 \({}^t\mathrm{H}\) 对于这些拓扑等度连续。

\begin{enumerate}
\def\labelenumi{\alph{enumi})}
\setcounter{enumi}{4}
\tightlist
\item
  证明：若 E 是桶型空间，则下列性质等价：
\end{enumerate}

α) H 在 \(\mathcal{L}(E; F)\) 中是简单有界的；

β) H 等度连续；

\(\gamma)\) \({}^{t}\mathrm{H}\) 在 \(\mathcal{L}(\mathrm{F}';\mathrm{E}')\) 中（当 \(\mathrm{E}'\) 赋以弱拓扑 \(\sigma (\mathrm{E}',\mathrm{E})\) 时）是简单有界的。

\(\delta)\) \({}^{t}\mathrm{H}\) 当 \(\mathrm{E}'\) 与 \(\mathrm{F}'\) 赋以强拓扑时等度连续。

\begin{enumerate}
\def\labelenumi{\alph{enumi})}
\setcounter{enumi}{5}
\tightlist
\item
  证明：若 E 是次桶型空间（III, 第 44 页, 习题 7），则性质 \(\beta\) 与 \(\delta\)（\(e\) ) 的）等价，且也与下列性质等价：
\end{enumerate}

ε) H 在 \(\mathcal{L}(\mathrm{E};\mathrm{F})\) 中对于有界收敛拓扑有界；

\(\phi)\) \({}^{t}\mathrm{H}\) 在 \(\mathcal{L}(\mathrm{F}';\mathrm{E}')\) 中（当 \(\mathrm{E}'\) 赋以强拓扑时）是简单有界的。

\begin{enumerate}
\def\labelenumi{\alph{enumi})}
\setcounter{enumi}{6}
\tightlist
\item
  证明：若 E 是拟完备的，则 e) 的性质 \(\alpha\) )、\(\gamma\) ) 与 \(\delta\) ) 等价。
\end{enumerate}

\begin{enumerate}
\def\labelenumi{\arabic{enumi})}
\setcounter{enumi}{19}
\tightlist
\item
  设 \(\mathbf{E}\) 是一个 Hausdorff 局部凸空间，\(\mathbf{E}'\) 是它的对偶。
\end{enumerate}

\begin{enumerate}
\def\labelenumi{\alph{enumi})}
\item
  设 M 是 E 的一个闭向量子空间。拓扑 \(\tau(\mathbf{M}, \mathbf{E}'/\mathbf{M}^{\circ})\) 与 \(\tau(\mathbf{E}, \mathbf{E}')\) 在 M 上诱导的拓扑相同，当且仅当 \(E'/M^{\circ}\) 中每个对于弱拓扑 \(\sigma(\mathrm{E}'/\mathrm{M}^{\circ}, \mathrm{M})\) 紧的凸均衡集，都是某个凸均衡且对于 \(\sigma(\mathrm{E}', \mathrm{E})\) 紧的子集（属于 \(E'\)）的典范像（参看 V, 第 73 页, 习题 15）。
\item
  设 N 是 E 的一个稠密向量子空间。若 E 的拓扑在 N 上诱导的拓扑与 \(\tau(\mathrm{N}, \mathrm{E}')\) 相同，证明 E 的拓扑与 \(\tau(\mathrm{E}, \mathrm{E}')\) 相同。
\end{enumerate}

\begin{enumerate}
\def\labelenumi{\arabic{enumi})}
\setcounter{enumi}{20}
\tightlist
\item
  \begin{enumerate}
  \def\labelenumii{\alph{enumii})}
  \tightlist
  \item
    设 E 是一个向量空间，\(E^{*}\) 是它的代数对偶。证明拓扑 \(\tau(\mathrm{E},\mathrm{E}^{*})\) 是最细的局部凸拓扑，且拓扑 \(\tau(\mathrm{E}^{*},\mathrm{E})\) 与 \(\sigma(\mathrm{E}^{*},\mathrm{E})\) 相同。b) 设 \(E^{**}\) 是 \(E^{*}\) 的代数对偶。证明：若 E 是无穷维的，则 E 在 \(E^{**}\) 中对于与 \(E^{**}\) 和 \(E^{*}\) 之间的对偶相容的所有拓扑都稠密，但 \(\tau(\mathrm{E}^{**},\mathrm{E}^{*})\) 在 E 上诱导的拓扑不同于 \(\tau(\mathrm{E},\mathrm{E}^{*})\)。
  \end{enumerate}
\item
  设 E 是一个 Hausdorff 局部凸空间，\(E'\) 是它的对偶，M 是 E 的一个闭子空间。a) 证明：若 E 中一个紧集的闭凸包是紧的，则 \(E'/M^{\circ}\)（视为 M 的对偶）上的紧收敛拓扑是 \(E'\) 上的紧收敛拓扑关于 \(M^{\circ}\) 的商拓扑。
\end{enumerate}

\begin{enumerate}
\def\labelenumi{\alph{enumi})}
\setcounter{enumi}{1}
\tightlist
\item
  证明：若 \(\mathbf{M}\) 是次桶型的，且 \(\mathrm{E}^{\prime} / \mathrm{M}^{\circ}\) 赋以拓扑 \(\beta (\mathrm{E}^{\prime} / \mathrm{M}^{\circ},\mathbf{M})\) 后是有界型的，则拓扑 \(\beta (\mathrm{E}^{\prime} / \mathrm{M}^{\circ},\mathbf{M})\) 是 \(\beta (\mathrm{E}^{\prime},\mathrm{E})\) 关于 \(\mathbf{M}^{\circ}\) 的商拓扑。
\end{enumerate}

\begin{enumerate}
\def\labelenumi{\arabic{enumi})}
\setcounter{enumi}{22}
\item
  给出一个 Hausdorff 局部凸空间族 \((\mathrm{E}_{i})_{i\in\mathrm{I}}\) 的例子，使得从 \(\bigoplus_{i\in I} E_{i}^{\prime}\) 到 \(P = \prod_{i\in I} E_{i}\) 的对偶上的典范映射不是从诸 \(E_{i}^{\prime}\)（赋以弱拓扑 \(\sigma(\mathrm{E}_{i}^{\prime}, \mathrm{E}_{i})\)）的拓扑直和到对偶 \(P^{\prime}\)（赋以 \(\sigma(\mathrm{P}^{\prime}, \mathrm{P})\)）上的同构。
\item
  设 \((\mathrm{E}_{\alpha})_{\alpha\in\mathrm{A}}\) 是一个局部凸 Hausdorff 空间族，E 是一个向量空间，且对每个 \(\alpha\in A\)，设 \(f_{\alpha}\) 是一个从 \(E_{\alpha}\) 到 E 中的线性映射。在 E 上考虑使诸函数 \(f_{\alpha}\) 连续的最细局部凸拓扑 T（II, 第 27 页）；假设 T 是 Hausdorff 的，并用 \(E_{\alpha}^{\prime}\) 表示 \(E_{\alpha}\) 的对偶，用 \(E^{\prime}\) 表示赋以 T 的 E 的对偶。证明：若对每个 \(\alpha\in A\)，\(E_{\alpha}\) 的拓扑与 \(\tau(\mathrm{E}_{\alpha},\mathrm{E}_{\alpha}^{\prime})\) 相同，则拓扑 T 与 \(\tau(\mathrm{E},\mathrm{E}^{\prime})\) 相同。
\item
  \begin{enumerate}
  \def\labelenumii{\alph{enumii})}
  \tightlist
  \item
    证明：每个实（相应地，复）Banach 空间都等距于形如 \(\mathcal{C}(\mathbf{S};\mathbf{R})\)（相应地，\(\mathcal{C}(\mathbf{S};\mathbf{C})\)）的 Banach 空间的一个闭向量子空间，这种空间由定义在一个紧空间 S 上的所有实（相应地，复）连续函数组成（GT, X, § 4, No. 1）（利用 IV, 第 7 页 的公式 (3)）。
  \end{enumerate}
\end{enumerate}

\begin{enumerate}
\def\labelenumi{\alph{enumi})}
\setcounter{enumi}{1}
\tightlist
\item
  由 a) 推出：每个 Hausdorff 局部凸空间 E 都同构于形如 \(\mathcal{C}_{c}(\mathrm{L};\mathbf{R})\)（相应地，\(\mathcal{C}_{c}(\mathrm{L};\mathbf{C})\)）的局部凸空间的一个子空间（GT, X, § 1, No. 6）。特别地，每个 Fréchet 空间都同构于一个空间 \(\mathcal{C}_{c}(\mathrm{L};\mathbf{R})\)（相应地，\(\mathcal{C}_{c}(\mathrm{L};\mathbf{C})\)）的闭子空间，其中 L 是局部紧且可分的。
\end{enumerate}
% semantic-fragments: legacy-input-seam
\relax{}
